Below is the revised write-up as a single LaTeX file, directly in chat. I kept the source-dependent claims tied to official documentation: RDKit’s stereo scope and APIs, the `FindPotentialStereo()` caveat, Daylight’s ring-closure lexical-order rule for chirality, and OpenSMILES ring/dot constraints. ([RDKit][1])

```latex
\documentclass[10pt,a4paper]{article}

\usepackage[margin=30mm]{geometry}
\usepackage[T1]{fontenc}
\usepackage[utf8]{inputenc}
\usepackage{amsmath,amssymb,amsthm,mathtools}
\usepackage{enumitem}
\usepackage{array}
\usepackage{booktabs}
\usepackage{xcolor}
\usepackage{listings}
\usepackage[pdftex,
            pdftitle={A Python implementation model for exact finite enumeration of bounded stereochemical SMILES},
            pdfauthor={},
            colorlinks,
            linkcolor=black,
            urlcolor=black,
            citecolor=black]{hyperref}

\allowdisplaybreaks
\emergencystretch=3em

\lstdefinestyle{pythonstyle}{
  language=Python,
  basicstyle=\ttfamily\small,
  columns=fullflexible,
  keepspaces=true,
  showstringspaces=false,
  breaklines=true,
  frame=single,
  framerule=0.2pt,
  rulecolor=\color{black!25},
  keywordstyle=\bfseries,
  commentstyle=\itshape\color{black!55}
}
\lstset{style=pythonstyle}

\newcommand{\calA}{\mathcal A}
\newcommand{\calB}{\mathcal B}
\newcommand{\calC}{\mathcal C}
\newcommand{\calD}{\mathcal D}
\newcommand{\calE}{\mathcal E}
\newcommand{\calG}{\mathcal G}
\newcommand{\calI}{\mathcal I}
\newcommand{\calK}{\mathcal K}
\newcommand{\calL}{\mathcal L}
\newcommand{\calM}{\mathcal M}
\newcommand{\calP}{\mathcal P}
\newcommand{\calR}{\mathcal R}
\newcommand{\calS}{\mathcal S}
\newcommand{\calT}{\mathcal T}
\newcommand{\calX}{\mathcal X}
\newcommand{\Bool}{\{0,1\}}
\newcommand{\None}{\bot}
\newcommand{\dotcup}{\mathbin{\dot\cup}}
\newcommand{\Sol}{\operatorname{Sol}}
\newcommand{\render}{\operatorname{render}}
\newcommand{\parse}{\operatorname{parse}}
\newcommand{\im}{\operatorname{im}}
\newcommand{\pos}{\operatorname{pos}}
\newcommand{\Decode}{\mathsf{Decode}}
\newcommand{\Tet}{\mathsf{Tet}}
\newcommand{\Dir}{\mathsf{Dir}}
\newcommand{\Pot}{\mathsf{Pot}}
\newcommand{\Ann}{\mathsf{Ann}}
\newcommand{\OK}{\mathsf{OK}}

\newtheorem{lemma}{Lemma}[section]
\newtheorem{proposition}[lemma]{Proposition}
\newtheorem{theorem}[lemma]{Theorem}
\newtheorem{corollary}[lemma]{Corollary}
\theoremstyle{definition}
\newtheorem{definition}[lemma]{Definition}
\newtheorem{assumption}[lemma]{Assumption}
\newtheorem{example}[lemma]{Example}
\theoremstyle{remark}
\newtheorem{remark}[lemma]{Remark}

\title{A Python implementation model for exact finite enumeration of a bounded stereochemical SMILES sublanguage}
\author{}
\date{\today}

\begin{document}
\maketitle

\begin{abstract}
This note gives an implementation-oriented model for exact support enumeration
of a bounded RDKit/OpenSMILES-like stereochemical SMILES sublanguage.  The
central software boundary is that an RDKit \texttt{Mol} is an input and audit
object only.  The source of truth for enumeration is an immutable finite object,
\texttt{MoleculeFacts}, containing atom, bond, component, stereochemical, and
ligand-occurrence facts.  Traversal skeletons, syntax slots, solver assignments,
annotation policies, and rendering consume only \texttt{MoleculeFacts} together
with explicit finite presentation and parser-semantics relations.  The result is
not a wrapper around a SMILES writer.  It is a finite enumerator whose satisfying
assignments render to exactly the declared support, up to the optional decision
to deduplicate equal rendered strings.
\end{abstract}

\noindent\textit{Keywords}: SMILES, stereochemistry, RDKit, Python, exact
enumeration, finite CSP, traversal grammar, tetrahedral chirality, directional
bonds, ring closures.

\section{Objective and invariant}
\label{s:objective}

The implementation should realize the finite map
\[
        (\calM,\calP,\calS)
        \longmapsto
        L(\calM,\calP,\calS)
        =
        \im\bigl(\render:\Sol C(\calM,\calP,\calS)\to \Sigma^*\bigr),
\]
where:
\begin{itemize}[leftmargin=*]
\item \(\calM\) is a finite molecular fact object;
\item \(\calP\) is a finite presentation policy;
\item \(\calS\) is a finite parser-semantics object for the chosen bounded
      dialect;
\item \(C(\calM,\calP,\calS)\) is the finite constraint object over traversal
      skeletons and syntax-slot assignments.
\end{itemize}

The rendered support is a set of strings.  A canonical-witness constraint may be
added for efficiency, but the mathematically clean support is the image of the
rendering map, not the set of internal witnesses.

The implementation should never ask RDKit whether a generated candidate is valid
during enumeration.  Such calls are appropriate for regression tests and external
audits, but not for accepting, rejecting, or repairing generated strings.
Correctness is by construction:
\begin{itemize}[leftmargin=*]
\item grammar-valid syntax comes from the traversal skeleton and syntax-slot
      domains;
\item graph identity comes from edge coverage, atom coverage, and decode
      constraints;
\item stereochemistry comes from explicit finite tetrahedral and directional
      relations;
\item absence of accidental stereo comes from constraints at all potential
      parser sites;
\item presentation conventions, including annotation maximality and ring-label
      normalization, are explicit policy predicates.
\end{itemize}

\section{The non-negotiable RDKit boundary}
\label{s:boundary}

The intended architecture is:
\[
\begin{array}{c}
\texttt{rdkit.Chem.Mol} \\
\downarrow\quad \text{one-way extraction} \\
\texttt{MoleculeFacts} \\
\downarrow\quad \text{skeletons, slots, constraints, rendering} \\
\text{support strings}.
\end{array}
\]

Only the adapter module may import RDKit.  The core package should be importable
and testable without RDKit installed.

\begin{lstlisting}
# allowed
from smiles_support.rdkit_adapter import molecule_facts_from_rdkit

facts = molecule_facts_from_rdkit(mol, extraction_policy)
texts = list(enumerate_support(facts, policy, semantics))

# forbidden outside rdkit_adapter.py and audit tests
from rdkit import Chem
Chem.MolFromSmiles(candidate)
Chem.MolToSmiles(mol)
\end{lstlisting}

The RDKit documentation describes support for tetrahedral atomic stereochemistry
and cis/trans stereochemistry at double bonds, while non-tetrahedral
stereochemistry remains outside the ordinary bounded subset considered here
\cite{rdkit-book-stereo}.  RDKit APIs such as \texttt{Bond.GetStereo()},
\texttt{Bond.GetStereoAtoms()}, and related atom/bond accessors are appropriate
at the adapter boundary, not inside the enumerator \cite{rdkit-rdchem}.

\section{Package topology}
\label{s:topology}

A practical package layout is:

\begin{lstlisting}
smiles_support/
  ids.py              # typed integer identifiers
  facts.py            # immutable MoleculeFacts; no RDKit
  graph_index.py      # derived adjacency/index caches; no RDKit
  policy.py           # finite presentation policy and text domains
  semantics.py        # finite parser relations; no RDKit
  skeleton.py         # traversal skeleton generation
  slots.py            # syntax-slot allocation from facts + skeleton
  constraints.py      # CSP predicates and small finite solvers
  annotation.py       # maximal/canonical annotation selectors
  ring_labels.py      # ring interval and label assignment constraints
  render.py           # deterministic renderer
  enumerate.py        # orchestration and deduplication
  rdkit_adapter.py    # the only module importing rdkit
  audit_rdkit.py      # optional external audit; not used by enumeration
\end{lstlisting}

The split among \texttt{facts.py}, \texttt{policy.py}, and
\texttt{semantics.py} is substantive:
\begin{itemize}[leftmargin=*]
\item \texttt{MoleculeFacts} says what finite labeled molecule is being
      enumerated.
\item \texttt{SmilesPolicy} says which finite presentation choices are admitted.
\item \texttt{ParserSemantics} says how the bounded dialect decodes those
      choices.
\end{itemize}
Conflating these layers is the main path back to writer-specific procedural
rules.

\section{Core finite data model}
\label{s:data-model}

\subsection{Typed identifiers}

Use typed integer identifiers.  They are stable, serializable, cheap, and make
it harder to confuse atoms, bonds, stereo sites, occurrences, and syntax slots.

\begin{lstlisting}
# ids.py
from typing import NewType

AtomId = NewType("AtomId", int)
BondId = NewType("BondId", int)
ComponentId = NewType("ComponentId", int)
SiteId = NewType("SiteId", int)
OccurrenceId = NewType("OccurrenceId", int)

AtomSlotId = NewType("AtomSlotId", int)
BondSlotId = NewType("BondSlotId", int)
CarrierSlotId = NewType("CarrierSlotId", int)
RingEndpointId = NewType("RingEndpointId", int)
EventPos = NewType("EventPos", int)
\end{lstlisting}

\subsection{Atom, bond, and component facts}

The graph facts are a frozen snapshot of an already perceived molecule.  They
are not instructions to re-perceive aromaticity, valence, hydrogens, or stereo.

\begin{lstlisting}
# facts.py
from __future__ import annotations

from dataclasses import dataclass
from enum import Enum

from .ids import AtomId, BondId, ComponentId, SiteId, OccurrenceId


class BondOrder(Enum):
    SINGLE = "single"
    DOUBLE = "double"
    TRIPLE = "triple"
    AROMATIC = "aromatic"


@dataclass(frozen=True, slots=True)
class AtomFacts:
    id: AtomId
    atomic_num: int
    symbol: str
    isotope: int | None
    formal_charge: int
    is_aromatic: bool
    explicit_h_count: int
    implicit_h_count: int
    no_implicit: bool


@dataclass(frozen=True, slots=True)
class BondFacts:
    id: BondId
    a: AtomId
    b: AtomId
    order: BondOrder
    is_aromatic: bool
    is_conjugated: bool


@dataclass(frozen=True, slots=True)
class ComponentFacts:
    id: ComponentId
    atoms: tuple[AtomId, ...]
    bonds: tuple[BondId, ...]
\end{lstlisting}

For the ordinary molecular graph considered here, assume at most one bond between
a pair of atoms.  If a broader multigraph dialect is admitted, the graph index
must store a multimap rather than a single \((a,b)\mapsto e\) lookup.

\subsection{Ligand occurrences}

A ligand occurrence is a semantic coordinate used by a stereo site.  It need not
be a graph vertex.  Examples include neighboring graph atoms, explicit hydrogen
atoms represented as graph vertices, implicit hydrogen occurrences, and any
policy-admitted pseudo-neighbor.

\begin{lstlisting}
class LigandKind(Enum):
    NEIGHBOR_ATOM = "neighbor_atom"
    IMPLICIT_H = "implicit_h"
    PSEUDO = "pseudo"


@dataclass(frozen=True, slots=True)
class LigandOccurrence:
    id: OccurrenceId
    site: SiteId
    kind: LigandKind
    atom: AtomId | None
    bond: BondId | None
    ordinal: int = 0
\end{lstlisting}

The occurrence object stores the finite ligand universe.  It does not store the
emitted SMILES ligand order.  For a tetrahedral center, that order is induced by
the traversal skeleton: root choice, parent edge, branch order, continuation
choice, and ring endpoint placement can all change it.

\subsection{Stereo facts}

The stereo layer should include both specified sites and potential but
unspecified sites.  Specified sites have a non-\texttt{NONE} target.  Potential
unspecified sites have target \texttt{NONE}; these are the loci for
no-accidental-stereo constraints.

\begin{lstlisting}
class SiteStatus(Enum):
    SPECIFIED = "specified"
    UNSPECIFIED = "unspecified"


class TetraValue(Enum):
    NONE = "none"
    PLUS = "plus"
    MINUS = "minus"


class DirectionalValue(Enum):
    NONE = "none"
    TOGETHER = "together"
    OPPOSITE = "opposite"


@dataclass(frozen=True, slots=True)
class TetrahedralSiteFacts:
    id: SiteId
    center: AtomId
    status: SiteStatus
    target: TetraValue
    ligand_occurrences: tuple[OccurrenceId, ...]
    reference_order: tuple[OccurrenceId, ...]


@dataclass(frozen=True, slots=True)
class DirectionalSiteFacts:
    id: SiteId
    center_bond: BondId
    left_endpoint: AtomId
    right_endpoint: AtomId
    status: SiteStatus
    target: DirectionalValue
    left_ligands: tuple[OccurrenceId, ...]
    right_ligands: tuple[OccurrenceId, ...]
    reference_pair: tuple[OccurrenceId, OccurrenceId] | None = None


@dataclass(frozen=True, slots=True)
class StereoFacts:
    tetrahedral: tuple[TetrahedralSiteFacts, ...]
    directional: tuple[DirectionalSiteFacts, ...]
\end{lstlisting}

The values \texttt{PLUS/MINUS} and \texttt{TOGETHER/OPPOSITE} are internal
coordinates.  They are not CIP labels.  The adapter normalizes RDKit or input
stereo into these coordinates using the stored reference order or reference
pair.

\subsection{Potential stereo sites}

\begin{definition}[Potential stereo site]
For a fixed fact object \(\calM\), policy \(\calP\), and parser semantics
\(\calS\), a potential stereo site is a parser locus \(q\) for which there exists
some legal traversal skeleton and syntax-slot assignment under \(\calP\) whose
semantic relation in \(\calS\) can return a specified non-\texttt{NONE}
stereochemical value at \(q\).  The set
\[
   \Pot(\calM,\calP,\calS)
\]
is therefore parser- and policy-defined.  Graph-theoretic recognition and RDKit
functions may be used by the adapter to help construct it, but they do not define
it.
\end{definition}

In code, it is acceptable to generate a conservative finite superset of potential
sites, provided the semantic relation for impossible sites always returns
\texttt{NONE}.  It is not acceptable to miss a parser-relevant potential site:
then accidental stereo can enter without a constraint.  RDKit's
\texttt{FindPotentialStereo()} is useful adapter input, but its documentation
describes it as somewhat experimental and warns that the API and results may
change \cite{rdkit-findpotential}.  That makes it a source of candidate facts,
not the mathematical definition.

\subsection{The fact object}

The final fact object should be immutable in meaning.  The snippets below use
tuples; a production implementation may add cached indexes outside the object.

\begin{lstlisting}
@dataclass(frozen=True, slots=True)
class MoleculeFacts:
    atoms: tuple[AtomFacts, ...]
    bonds: tuple[BondFacts, ...]
    components: tuple[ComponentFacts, ...]
    stereo: StereoFacts
    ligand_occurrences: tuple[LigandOccurrence, ...]

    def validate(self) -> None:
        """Check id uniqueness, endpoint validity, component partition,
        stereo occurrence validity, and reference-order consistency."""
        ...
\end{lstlisting}

Avoid storing mutable dictionaries directly in frozen dataclasses unless they are
explicitly wrapped or treated as internal caches.  Prefer sorted tuples in fact
objects and derived maps in indexes.

\begin{lstlisting}
@dataclass(slots=True)
class GraphIndex:
    atom_by_id: dict[AtomId, AtomFacts]
    bond_by_id: dict[BondId, BondFacts]
    incident_bonds: dict[AtomId, tuple[BondId, ...]]
    neighbors: dict[AtomId, tuple[AtomId, ...]]
    bond_between: dict[tuple[AtomId, AtomId], BondId]
\end{lstlisting}

The index is derived data.  It should not perform chemistry perception.

\section{Policy and finite syntax domains}
\label{s:policy}

The policy supplies finite presentation domains.  It must be explicit about what
is excluded.  In particular, atom maps, CXSMILES fields, arbitrary numeric
extensions, arbitrary whitespace, unbounded ring labels, and writer-specific
canonicalization side effects should be absent or separately bounded.

\begin{lstlisting}
# policy.py
from dataclasses import dataclass
from enum import Enum
from typing import Mapping

from .ids import AtomId, BondId
from .facts import MoleculeFacts


class TetraToken(Enum):
    NONE = ""
    AT = "@"
    ATAT = "@@"


class DirectionMark(Enum):
    ABSENT = 0
    FWD = 1
    REV = -1


@dataclass(frozen=True, slots=True)
class RingLabel:
    value: int

    def text(self) -> str:
        if 1 <= self.value <= 9:
            return str(self.value)
        return f"%{self.value:02d}"


@dataclass(frozen=True, slots=True)
class AtomTextChoice:
    name: str
    # In production use an immutable tuple/map.
    text_by_tetra: Mapping[TetraToken, str]

    def permits(self, token: TetraToken) -> bool:
        return token in self.text_by_tetra

    def render(self, token: TetraToken) -> str:
        return self.text_by_tetra[token]


@dataclass(frozen=True, slots=True)
class BondTextChoice:
    name: str
    base_text: str
    permits_direction: bool


class AnnotationMode(Enum):
    HARD = "hard"
    SUPPORT_MAXIMAL = "support_maximal"
    CARDINALITY_MAXIMAL = "cardinality_maximal"
    CANONICAL = "canonical"


@dataclass(frozen=True, slots=True)
class SmilesPolicy:
    ring_labels: tuple[RingLabel, ...]
    annotation_mode: AnnotationMode
    deduplicate_rendered_strings: bool = True
    least_free_ring_labels: bool = True

    def atom_text_domain(
        self,
        facts: MoleculeFacts,
        atom: AtomId,
    ) -> tuple[AtomTextChoice, ...]:
        ...

    def bond_text_domain(
        self,
        facts: MoleculeFacts,
        bond: BondId,
        *,
        slot_kind: str,
    ) -> tuple[BondTextChoice, ...]:
        ...
\end{lstlisting}

Omission of a directional marker should be represented as the first-class domain
value \texttt{DirectionMark.ABSENT}.  It should not be represented by absence of
a variable.  Treating absence as a value makes no-accidental-stereo and
maximality predicates total finite relations.

\begin{remark}[Contextual atom and bond decoding]
For a bracketed atom spelling, decoding is often a unary relation between the
text and \texttt{AtomFacts}.  For unbracketed atoms, elided bonds, aromatic
spellings, and implicit hydrogen interpretations, decoding may depend on the
incident bond-order context.  The clean implementation is therefore not always a
unary atom table but a finite star relation
\[
   \Decode^{\mathrm{atom}}_v
      \bigl(a_v,\tau_v,(b_s,\delta_s)_{s\in\partial v}\bigr).
\]
Likewise, a ring-closure edge may require a joint endpoint relation
\[
   \Decode^{\mathrm{ring}}_e
      \bigl(b_{e,1},\delta_{e,1},b_{e,2},\delta_{e,2},\lambda_e\bigr),
\]
rather than two independent endpoint checks.  This remains finite.  It is better
to model the relation explicitly than to rely on a parser to repair or reject
candidate strings after generation.
\end{remark}

\section{Traversal skeletons}
\label{s:skeletons}

\subsection{Canonical skeleton object}

The minimal clean traversal witness is a rooted spanning forest plus an ordered
list of local post-atom events.  For each connected component, the policy chooses
or enumerates a root.  Each non-root atom has one parent.  Tree edges are emitted
as traversal bonds.  Non-tree edges are emitted as paired ring endpoints.  At
each atom, local events are ordered; each child event is either a branch child or,
if it is the final child event, a continuation child.

\begin{lstlisting}
# skeleton.py
from dataclasses import dataclass
from enum import Enum

from .ids import AtomId, BondId


class ChildRole(Enum):
    BRANCH = "branch"
    CONTINUATION = "continuation"


@dataclass(frozen=True, slots=True)
class ChildEvent:
    bond: BondId
    parent: AtomId
    child: AtomId
    role: ChildRole


@dataclass(frozen=True, slots=True)
class RingEndpointEvent:
    bond: BondId
    at: AtomId


LocalEvent = ChildEvent | RingEndpointEvent


@dataclass(frozen=True, slots=True)
class TraversalSkeleton:
    roots: tuple[AtomId, ...]
    parent: dict[AtomId, AtomId | None]
    tree_bonds: frozenset[BondId]
    ring_bonds: frozenset[BondId]
    events_at: dict[AtomId, tuple[LocalEvent, ...]]
\end{lstlisting}

This object is equivalent to an ordered DFS event tree with ring back-edges.  It
is more useful than a raw character string because tetrahedral local order and
directional carrier scopes can be computed before rendering.

\subsection{Branch/continuation as linearization}

Branch/continuation choice is part of the presentation skeleton.  It can be
viewed either as a child-role variable or as a linearization constraint on DFS
events:
\[
   \text{at most one child traversal is unparenthesized, and if present it is last.}
\]
The two views are equivalent.  In code, the child-role view is easier to render
and debug.

\subsection{Skeleton generation}

The first implementation should use an explicit generator rather than a
monolithic SAT encoding of traversal.  This is still an exact finite
decomposition of the CSP; it is not a writer repair procedure.

\begin{lstlisting}
def enumerate_skeletons(
    facts: MoleculeFacts,
    index: GraphIndex,
    policy: SmilesPolicy,
) -> Iterator[TraversalSkeleton]:
    for roots in enumerate_component_roots(facts, policy):
        for tree_bonds in enumerate_spanning_forests(facts, index, roots):
            parent = orient_forest(index, roots, tree_bonds)
            ring_bonds = frozenset(b.id for b in facts.bonds) - frozenset(tree_bonds)

            local_domains: list[tuple[AtomId, tuple[tuple[LocalEvent, ...], ...]]] = []

            for atom in atom_ids(facts):
                children = children_of(parent, atom, index)
                ring_eps = tuple(
                    RingEndpointEvent(bond=bid, at=atom)
                    for bid in ring_bonds
                    if is_incident(index, bid, atom)
                )
                local_domains.append(
                    (atom, tuple(local_event_orders(children, ring_eps)))
                )

            for events_at in product_local_orders(local_domains):
                yield TraversalSkeleton(
                    roots=roots,
                    parent=parent,
                    tree_bonds=frozenset(tree_bonds),
                    ring_bonds=ring_bonds,
                    events_at=events_at,
                )
\end{lstlisting}

For large graphs the generator will be enormous.  That is not an implementation
bug: exact support enumeration is output-sensitive and generally exponential.
The implementation should still prune early using fragment order, root policy,
maximum simultaneously open ring labels, atom-text impossibility, and stereo
infeasibility.

\section{Syntax slots}
\label{s:slots}

A skeleton determines a finite set of syntax slots.  Assignments choose values
for these slots.

\begin{lstlisting}
# slots.py
from dataclasses import dataclass
from enum import Enum

from .ids import AtomId, BondId, SiteId, BondSlotId, CarrierSlotId, RingEndpointId


class BondSlotKind(Enum):
    TREE = "tree"
    RING_ENDPOINT = "ring_endpoint"


@dataclass(frozen=True, slots=True)
class AtomSlot:
    atom: AtomId


@dataclass(frozen=True, slots=True)
class BondSlot:
    id: BondSlotId
    bond: BondId
    kind: BondSlotKind
    written_from: AtomId
    written_to: AtomId | None
    ring_endpoint: RingEndpointId | None = None


@dataclass(frozen=True, slots=True)
class CarrierSlot:
    id: CarrierSlotId
    bond_slot: BondSlotId
    bond: BondId
    written_from: AtomId
    written_to: AtomId | None


@dataclass(frozen=True, slots=True)
class SlotBundle:
    atom_slots: tuple[AtomSlot, ...]
    bond_slots: tuple[BondSlot, ...]
    carrier_slots: tuple[CarrierSlot, ...]
    ring_endpoints: tuple[RingEndpointId, ...]
    directional_scope: dict[SiteId, tuple[CarrierSlotId, ...]]
\end{lstlisting}

A directional carrier slot must remember the direction in which the bond text is
written.  A slash/backslash character has no global geometric sign.  Its
site-relative meaning depends on the written direction of the carrier and on the
stereo site currently being interpreted.

\section{Parser semantics as finite relations}
\label{s:semantics}

The semantics object is the implementation counterpart of the mathematical
relations \(\Decode\), \(\Tet\), and \(\Dir\).  It should be pure and finite:
for a fixed fact object, skeleton, slot bundle, and assignment, it returns a
value from a finite set or an invalid marker.

\begin{lstlisting}
# semantics.py
from typing import Mapping, Protocol

from .ids import AtomId, BondId, SiteId, OccurrenceId, CarrierSlotId
from .facts import MoleculeFacts
from .policy import AtomTextChoice, BondTextChoice, TetraToken, DirectionMark
from .skeleton import TraversalSkeleton
from .slots import SlotBundle


class Invalid:
    pass


INVALID = Invalid()


class ParserSemantics(Protocol):
    def atom_decode_ok(
        self,
        facts: MoleculeFacts,
        atom: AtomId,
        atom_text: AtomTextChoice,
        tetra_token: TetraToken,
        incident_bond_texts: tuple[BondTextChoice, ...],
    ) -> bool:
        ...

    def bond_decode_ok(
        self,
        facts: MoleculeFacts,
        bond: BondId,
        bond_text: BondTextChoice,
        direction_mark: DirectionMark,
    ) -> bool:
        ...

    def ring_pair_decode_ok(
        self,
        facts: MoleculeFacts,
        bond: BondId,
        endpoint_1: BondTextChoice,
        mark_1: DirectionMark,
        endpoint_2: BondTextChoice,
        mark_2: DirectionMark,
    ) -> bool:
        ...

    def local_tetra_order(
        self,
        facts: MoleculeFacts,
        skel: TraversalSkeleton,
        slots: SlotBundle,
        site: SiteId,
    ) -> tuple[OccurrenceId, ...]:
        ...

    def tetra_value(
        self,
        facts: MoleculeFacts,
        site: SiteId,
        local_order: tuple[OccurrenceId, ...],
        token: TetraToken,
    ) -> TetraValue | Invalid:
        ...

    def directional_value(
        self,
        facts: MoleculeFacts,
        skel: TraversalSkeleton,
        slots: SlotBundle,
        site: SiteId,
        marks: Mapping[CarrierSlotId, DirectionMark],
    ) -> DirectionalValue | Invalid:
        ...
\end{lstlisting}

For ordinary alkene/imine/oxime-style directional stereo, a signed parity CSP is
often the efficient implementation.  The exposed primitive should still be the
finite relation \texttt{directional\_value}.  It can internally use parity
tables, but it can also handle endpoint degeneracy, ring endpoint orientation,
missing-carrier patterns, and contradictory marker patterns without changing the
rest of the enumerator.

\subsection{Tetrahedral constraint}

For each tetrahedral site \(q\), the skeleton induces a local occurrence order
\(L_q\).  The tetrahedral token at the center is constrained by
\[
   \Tet_q(L_q,\tau_q)=t_q,
\]
where \(t_q\) is \texttt{PLUS/MINUS} for specified sites and \texttt{NONE} for
potential but unspecified sites.  The relation must include the dialect-specific
rule for implicit hydrogen placement and for ring-closure occurrence ordering.
Daylight's SMILES theory states that, for chirality, the order of a ring-closure
bond is determined by the lexical position of the ring digit on the chiral atom,
not by the later position of the other endpoint \cite{daylight-smiles}.

\subsection{Directional constraint}

For each directional site \(q\), compute its carrier scope
\[
   K_q\subseteq \texttt{CarrierSlotId}.
\]
The constraint is
\[
   \Dir_q\bigl((\delta_k)_{k\in K_q}\bigr)=d_q,
\]
where \(d_q\) is \texttt{TOGETHER/OPPOSITE} for specified sites and
\texttt{NONE} for potential but unspecified sites.

A single carrier slot may belong to more than one site scope.  Therefore
no-accidental-stereo is represented by local finite relations placed in one
global CSP, not by independently solving each site.

\begin{lstlisting}
def check_tetrahedral_constraints(
    facts: MoleculeFacts,
    skel: TraversalSkeleton,
    slots: SlotBundle,
    assignment: "FullAssignment",
    semantics: ParserSemantics,
) -> bool:
    for site in facts.stereo.tetrahedral:
        token = assignment.tetra_token.get(site.center, TetraToken.NONE)
        order = semantics.local_tetra_order(facts, skel, slots, site.id)
        value = semantics.tetra_value(facts, site.id, order, token)
        if value is INVALID or value != site.target:
            return False
    return True


def check_directional_constraints(
    facts: MoleculeFacts,
    skel: TraversalSkeleton,
    slots: SlotBundle,
    assignment: "FullAssignment",
    semantics: ParserSemantics,
) -> bool:
    for site in facts.stereo.directional:
        scope = slots.directional_scope[site.id]
        marks = {
            k: assignment.direction_mark.get(k, DirectionMark.ABSENT)
            for k in scope
        }
        value = semantics.directional_value(
            facts=facts,
            skel=skel,
            slots=slots,
            site=site.id,
            marks=marks,
        )
        if value is INVALID or value != site.target:
            return False
    return True
\end{lstlisting}

The same functions enforce preservation of specified stereo and absence of
accidental stereo.  Only the stored target differs.

\begin{remark}[Local relations versus global constraints]
If the bounded parser semantics assigns each site value from a bounded local
scope, then no-accidental-stereo is implemented by one finite relation per site,
with shared variables creating a global CSP.  If a chosen parser semantics
contains a genuinely global cleanup or conflict-resolution rule, that rule must
be represented as a finite global relation or excluded from the semantics.  The
important point is not locality at all costs; it is that every parser-relevant
effect is represented before rendering rather than delegated to post-hoc
validation.
\end{remark}

\section{Ring labels}
\label{s:ring-labels}

Each non-tree edge produces two ring endpoint events.  A ring-label assignment
is a function from ring edges to a finite label set.  The basic correctness
conditions are:
\begin{enumerate}[leftmargin=*]
\item each ring edge has exactly two endpoint events;
\item both endpoints of the same ring edge use the same label;
\item a label is not reused while it is live;
\item endpoint bond text and directional markers jointly decode to the fixed
      graph bond;
\item optional policy constraints, such as least-currently-free label, are
      applied as finite predicates;
\item grammar-specific exclusions, such as forbidding a dot immediately before a
      ring-closure digit, are encoded in the ring endpoint grammar rather than
      checked after parsing.
\end{enumerate}

The interval condition for reusable labels is:
\[
  \lambda_e=\lambda_f,\quad \pos_1(e)<\pos_1(f)
  \quad\Longrightarrow\quad
  \pos_2(e)<\pos_1(f),
\]
where \(\pos_1(e)<\pos_2(e)\) are the two event positions of edge \(e\).  Give
every endpoint event a strict total event position, not merely an atom position,
so that multiple ring endpoints on the same atom are unambiguous.

\begin{lstlisting}
def legal_ring_labels(
    skel: TraversalSkeleton,
    policy: SmilesPolicy,
) -> Iterator[dict[BondId, RingLabel]]:
    intervals = compute_ring_intervals(skel)  # BondId -> (open_event_pos, close_event_pos)
    ring_edges = tuple(sorted(skel.ring_bonds, key=lambda e: intervals[e][0]))
    labels = policy.ring_labels
    out: dict[BondId, RingLabel] = {}

    def active_at(pos: EventPos) -> set[RingLabel]:
        active: set[RingLabel] = set()
        for e, lab in out.items():
            lo, hi = intervals[e]
            if lo < pos < hi:
                active.add(lab)
        return active

    def rec(i: int) -> Iterator[dict[BondId, RingLabel]]:
        if i == len(ring_edges):
            yield dict(out)
            return

        e = ring_edges[i]
        lo, _ = intervals[e]
        candidates = tuple(lab for lab in labels if lab not in active_at(lo))

        if policy.least_free_ring_labels:
            if not candidates:
                return
            candidates = (min(candidates, key=lambda x: x.value),)

        for lab in candidates:
            out[e] = lab
            yield from rec(i + 1)
            del out[e]

    yield from rec(0)
\end{lstlisting}

Interval non-overlap is necessary for reusable labels, but it is not the whole
ring-closure model.  Pairing, endpoint identity, joint bond decode, event
position, and grammar exclusions are also part of correctness.

\section{Assignments}
\label{s:assignments}

Assignments should be explicit typed objects.  A satisfying assignment is a
complete finite witness.  Rendering should only happen after all constraints and
policy predicates have been applied.

\begin{lstlisting}
# constraints.py
from dataclasses import dataclass

from .ids import AtomId, BondId, BondSlotId, CarrierSlotId
from .policy import AtomTextChoice, BondTextChoice, RingLabel, TetraToken, DirectionMark


@dataclass(frozen=True, slots=True)
class BaseAssignment:
    atom_text: dict[AtomId, AtomTextChoice]
    tetra_token: dict[AtomId, TetraToken]
    bond_text: dict[BondSlotId, BondTextChoice]
    ring_label: dict[BondId, RingLabel]


@dataclass(frozen=True, slots=True)
class FullAssignment:
    atom_text: dict[AtomId, AtomTextChoice]
    tetra_token: dict[AtomId, TetraToken]
    bond_text: dict[BondSlotId, BondTextChoice]
    ring_label: dict[BondId, RingLabel]
    direction_mark: dict[CarrierSlotId, DirectionMark]

    @classmethod
    def from_base(
        cls,
        base: BaseAssignment,
        direction_mark: dict[CarrierSlotId, DirectionMark],
    ) -> "FullAssignment":
        return cls(
            atom_text=base.atom_text,
            tetra_token=base.tetra_token,
            bond_text=base.bond_text,
            ring_label=base.ring_label,
            direction_mark=direction_mark,
        )
\end{lstlisting}

For efficiency, build assignments in stages.  Ring labels, atom text,
non-directional bond text, and tetrahedral syntax can often be checked before
enumerating slash/backslash markers.

\section{Annotation policy}
\label{s:annotation}

An annotation policy is best implemented as a selector over assignments that
already satisfy the semantic constraints.  It may be:
\begin{itemize}[leftmargin=*]
\item a hard predicate on complete assignments;
\item an optimization or maximality layer;
\item a canonicalization rule;
\item a pre-solver restriction of allowed slots.
\end{itemize}
These are not equivalent in general.  A pre-solver generator is equivalent to a
hard predicate only when it exactly characterizes the same complete-assignment
set.  It is equivalent to maximality only under additional monotonicity
assumptions that often fail for shared directional carriers.

\begin{lstlisting}
def marker_support(a: FullAssignment) -> frozenset[CarrierSlotId]:
    return frozenset(
        k for k, v in a.direction_mark.items()
        if v != DirectionMark.ABSENT
    )


def select_by_annotation_policy(
    assignments: list[FullAssignment],
    policy: SmilesPolicy,
) -> list[FullAssignment]:
    if policy.annotation_mode is AnnotationMode.HARD:
        return assignments

    supports = [marker_support(a) for a in assignments]

    if policy.annotation_mode is AnnotationMode.SUPPORT_MAXIMAL:
        return [
            a for a, s in zip(assignments, supports)
            if not any(s < t for t in supports)
        ]

    if policy.annotation_mode is AnnotationMode.CARDINALITY_MAXIMAL:
        if not assignments:
            return []
        m = max(len(s) for s in supports)
        return [a for a, s in zip(assignments, supports) if len(s) == m]

    if policy.annotation_mode is AnnotationMode.CANONICAL:
        if not assignments:
            return []
        return [min(assignments, key=canonical_annotation_key)]

    raise ValueError(policy.annotation_mode)
\end{lstlisting}

Support-wise maximality is the direct interpretation of:
\[
   \text{emit every eligible carrier marker that can participate in a surviving assignment}.
\]
It can yield multiple incomparable maximal supports.  Cardinality maximality
removes some but not all multiplicity.  Canonical maximality gives one witness by
adding a presentation convention.  The implementation should not hide this
choice.

For a pure SAT implementation, support maximality can be encoded by blocking all
strictly larger feasible supports.  This is finite but can be exponentially large
in the number of carrier slots.  An optimization or MaxSAT-style layer is the
honest formulation.  A backtracking implementation can apply the same selector
after enumerating feasible directional assignments for a fixed skeleton and base
assignment.

\section{Rendering}
\label{s:rendering}

The renderer is deterministic string concatenation.  It should not ask chemistry
questions.  It consumes a skeleton, slots, and a satisfying assignment.

\begin{lstlisting}
# render.py
def render(
    facts: MoleculeFacts,
    skel: TraversalSkeleton,
    slots: SlotBundle,
    a: FullAssignment,
) -> str:
    return ".".join(
        render_atom_subtree(root, facts, skel, slots, a)
        for root in skel.roots
    )


def render_atom_subtree(
    atom: AtomId,
    facts: MoleculeFacts,
    skel: TraversalSkeleton,
    slots: SlotBundle,
    a: FullAssignment,
) -> str:
    out: list[str] = [
        a.atom_text[atom].render(a.tetra_token.get(atom, TetraToken.NONE))
    ]

    for event in skel.events_at[atom]:
        if isinstance(event, RingEndpointEvent):
            bs = bond_slot_for_ring_endpoint(slots, event)
            out.append(render_bond_slot(bs, a))
            out.append(a.ring_label[event.bond].text())

        elif isinstance(event, ChildEvent):
            bs = bond_slot_for_tree_edge(slots, event.bond)
            child_text = (
                render_bond_slot(bs, a)
                + render_atom_subtree(event.child, facts, skel, slots, a)
            )
            if event.role is ChildRole.BRANCH:
                out.append("(" + child_text + ")")
            else:
                out.append(child_text)

        else:
            raise TypeError(event)

    return "".join(out)
\end{lstlisting}

The slash/backslash character is rendered by a finite table from
\texttt{DirectionMark} and slot orientation.  The same orientation convention
must be used by \texttt{ParserSemantics.directional\_value}.  This shared
convention should be tested directly with small hand-built examples.

\section{Enumeration pipeline}
\label{s:pipeline}

The top-level enumerator is a sequence of finite generators and predicates.

\begin{lstlisting}
# enumerate.py
def enumerate_support(
    facts: MoleculeFacts,
    policy: SmilesPolicy,
    semantics: ParserSemantics,
) -> Iterator[str]:
    facts.validate()
    index = build_graph_index(facts)
    seen: set[str] = set()

    for skel in enumerate_skeletons(facts, index, policy):
        slots = allocate_slots(facts, index, skel, policy, semantics)

        for base in enumerate_base_assignments(
            facts, index, skel, slots, policy, semantics
        ):
            if not base_constraints_ok(
                facts, index, skel, slots, base, policy, semantics
            ):
                continue

            feasible: list[FullAssignment] = []

            for full in enumerate_direction_assignments(
                facts, skel, slots, base, policy, semantics
            ):
                if all_constraints_ok(
                    facts, index, skel, slots, full, policy, semantics
                ):
                    feasible.append(full)

            for full in select_by_annotation_policy(feasible, policy):
                text = render(facts, skel, slots, full)

                if policy.deduplicate_rendered_strings:
                    if text in seen:
                        continue
                    seen.add(text)

                yield text
\end{lstlisting}

The \texttt{seen} set implements the image-of-render definition.  It removes
duplicate witnesses caused by automorphisms, symmetric roots, or distinct
presentation witnesses that collapse to the same string.  It is not parse
filtering.  If memory is a concern, replace it with a canonical-witness
constraint, a streaming external sort, or a hash-backed database.

\section{RDKit adapter}
\label{s:adapter}

The adapter snapshots RDKit facts into \texttt{MoleculeFacts}.  It should be
small, isolated, and version-tested.  It may use RDKit's perception functions,
but their results must be stored as explicit facts before enumeration starts.

\begin{lstlisting}
# rdkit_adapter.py
from rdkit import Chem

from .facts import MoleculeFacts


def molecule_facts_from_rdkit(
    mol: Chem.Mol,
    extraction_policy: "ExtractionPolicy",
) -> MoleculeFacts:
    atoms = tuple(extract_atoms(mol))
    bonds = tuple(extract_bonds(mol))
    components = tuple(extract_components(mol, atoms, bonds))

    stereo, occurrences = extract_stereo_and_occurrences(
        mol=mol,
        atoms=atoms,
        bonds=bonds,
        extraction_policy=extraction_policy,
    )

    facts = MoleculeFacts(
        atoms=atoms,
        bonds=bonds,
        components=components,
        stereo=stereo,
        ligand_occurrences=occurrences,
    )
    facts.validate()
    return facts
\end{lstlisting}

\subsection{Adapter responsibilities}

The adapter has four responsibilities.

\begin{enumerate}[leftmargin=*]
\item Copy fixed atom and bond attributes: element, isotope, charge, aromaticity,
      explicit and implicit hydrogens, bond order, and conjugation flags.
\item Normalize specified tetrahedral and directional stereo into internal site
      facts with explicit reference orders or reference pairs.
\item Construct the parser/policy-relative potential-site universe.  RDKit may
      propose candidates, but the adapter must reconcile those candidates with
      the bounded dialect.
\item Reject inputs outside the declared scope, such as enhanced stereo groups,
      non-tetrahedral atom stereo, atropisomerism, unknown stereo modes not
      represented by the model, or unsupported bond stereo classes.
\end{enumerate}

Do not store RDKit's chiral tag alone as a fact.  A chiral tag has meaning only
relative to an ordering convention.  Store the normalized tetrahedral value plus
the occurrence reference order that was used to compute it.

\subsection{Audit use of RDKit}

RDKit may be used after enumeration in tests:

\begin{lstlisting}
def audit_generated_strings_with_rdkit(
    mol: Chem.Mol,
    strings: Iterable[str],
) -> None:
    for smi in strings:
        parsed = Chem.MolFromSmiles(smi)
        assert parsed is not None
        # Compare against independent invariants or a new adapter snapshot.
        # Never call this from enumerate_support().
\end{lstlisting}

This audit is valuable for finding implementation mistakes, but it is not part
of the support definition.

\section{Solver architecture}
\label{s:solver}

The simplest correct implementation is a hybrid generator/CSP checker:
\begin{enumerate}[leftmargin=*]
\item enumerate traversal skeletons;
\item allocate slots;
\item enumerate ring labels and base text assignments with early decode checks;
\item solve directional marker assignments for each base witness;
\item apply the annotation policy;
\item render and deduplicate strings.
\end{enumerate}

This should be the first implementation because it exposes every semantic layer
and is easy to debug.  A later backend can compile parts to SAT or CP.  Good
first SAT targets are ring labels, directional markers, and maximal annotation.
The traversal skeleton is usually the least pleasant part to encode first,
because spanning-tree, DFS-order, branch, and ring-endpoint variables create a
large amount of symmetry.

\begin{lstlisting}
class DirectionSolver(Protocol):
    def solve(
        self,
        facts: MoleculeFacts,
        skel: TraversalSkeleton,
        slots: SlotBundle,
        base: BaseAssignment,
        policy: SmilesPolicy,
        semantics: ParserSemantics,
    ) -> Iterator[FullAssignment]:
        ...
\end{lstlisting}

The initial implementation can be recursive backtracking with pruning.  A later
implementation can use a SAT solver and return the same \texttt{FullAssignment}
objects.  Do not let solver-specific variable encodings leak into the domain
model.

\section{Alternative formal views}
\label{s:formalism}

The implementation model is naturally described as an attributed graph traversal
grammar plus a finite-domain CSP.  Equivalent or related formal views are
possible:
\begin{itemize}[leftmargin=*]
\item a regular tree grammar with finite attributes and side constraints;
\item a constraint automaton over traversal events;
\item an MSO-definable structure over a graph equipped with a traversal order;
\item a finite graph-to-string transducer with regular constraints;
\item a SAT/CP instance compiled from skeleton, slot, and policy variables.
\end{itemize}

These formalisms are useful for proofs or compilation.  For Python, the most
maintainable representation is still typed domain objects plus finite relations.
The critical point is that the entire bounded dialect is finite and explicit;
the exact implementation need not be monolithic.

\section{Correctness invariant}
\label{s:correctness}

\begin{theorem}[Implementation-level soundness]
Assume the adapter has produced a valid \texttt{MoleculeFacts} object, the policy
has finite domains, and the semantics object exactly represents the bounded
parser relations.  If \texttt{enumerate\_support} yields a string \(w\), then
\(w\) is in the declared finite SMILES sublanguage and parses, under the declared
semantics, to a graph isomorphic to the input fact graph with the specified
stereochemistry and with no additional potential-site stereo.
\end{theorem}

\begin{proof}
The skeleton generator emits only legal atom, branch, continuation,
disconnection, and ring-endpoint event structures.  The graph constraints ensure
that every atom is emitted once and every graph bond is emitted either as one
traversal bond or as one paired ring closure.  The decode constraints ensure
that atom and bond spellings decode to the stored atom and bond facts.  The ring
constraints ensure endpoint pairing, legal bounded label reuse, and joint
ring-edge decode.  The tetrahedral and directional relations match specified-site
targets and force \texttt{NONE} at potential unspecified sites.  The renderer
only concatenates texts chosen by a satisfying assignment.  Therefore the
rendered string has the stated parse-equivalence property.
\end{proof}

\begin{theorem}[Implementation-level completeness]
Under the same assumptions, for every string \(w\in L(\calM,\calP,\calS)\) in
the declared support, there is at least one satisfying assignment rendered by
the implementation as \(w\), unless an optional canonical-witness rule was chosen
that intentionally selects a representative subset of witnesses while preserving
rendered support.
\end{theorem}

\begin{proof}
Read the root choices, parent relation, branch/continuation choices, ring
endpoint positions, ring labels, atom text choices, bond text choices,
tetrahedral tokens, and directional marks directly from \(w\).  Because \(w\) is
in the declared sublanguage and parses to the input fact graph under the declared
semantics, these choices satisfy the corresponding finite constraints.  The
renderer reconstructs \(w\) from those choices.  If only rendered strings are
deduplicated, the string is still yielded once.  If a canonical-witness rule is
used, completeness is with respect to the canonicalized witness set by design.
\end{proof}

The theorem does not require an explicit quotient by graph automorphisms.  The
image-of-render formulation already quotients by equality of strings.  Graph
automorphisms matter for performance and duplicate witnesses, not for semantic
correctness.

\begin{remark}[Parser determinism]
The theorem assumes that the declared semantics assigns a well-defined finite
semantic value, or an invalid value, to each relevant site under a complete slot
assignment.  If the chosen parser semantics is nondeterministic, replace each
site value by a finite set of possible values and state whether support requires
all parses or at least one parse to match the fact graph.  The deterministic
relation is the cleaner implementation target.
\end{remark}

\section{Complexity and output sensitivity}
\label{s:complexity}

There are three different enumeration sizes.

\begin{center}
\begin{tabular}{>{\raggedright\arraybackslash}p{0.25\linewidth}
                >{\raggedright\arraybackslash}p{0.65\linewidth}}
\toprule
Quantity & Meaning \\
\midrule
satisfying assignments &
all traversal, slot, and marker witnesses; may contain many duplicates by
rendering \\
canonical witnesses &
one chosen witness per presentation equivalence class, if a canonical rule is
imposed \\
rendered support strings &
the actual language support, \(\im(\render)\) \\
\bottomrule
\end{tabular}
\end{center}

State complexity claims relative to one of these quantities.  The honest default
is output-sensitive enumeration of rendered support strings, with potentially
exponential preprocessing and delay in the size of the graph.  If the API yields
assignments instead of strings, say so.  If a canonical-witness constraint is
active, state the canonicalization policy because it becomes part of the
enumerated object.

\section{Sources of infinity to exclude or bound}
\label{s:infinity}

For a fixed finite graph, support remains finite only if all presentation domains
are finite.  Ring labels are the obvious source of infinity, but not the only
one.  The policy should also bound or exclude:
\begin{enumerate}[leftmargin=*]
\item arbitrary ring-label extensions;
\item arbitrary isotope, charge, hydrogen-count, or atom-map numeric spellings;
\item CXSMILES or name fields after the core SMILES;
\item redundant whitespace or comments, if the chosen parser accepts them;
\item unbounded alternative aromatic/Kekule or valence-normalization spellings;
\item wildcard, query, or SMARTS-like syntax if not part of the dialect.
\end{enumerate}

If any of these are admitted, their finite domains must be placed in the policy
object and in the decode relations.

\section{Known implementation pitfalls}
\label{s:pitfalls}

\begin{enumerate}[leftmargin=*]
\item \textbf{Slash orientation.}  A slash/backslash token is not a global sign.
      Its meaning depends on the written direction of the carrier slot and on
      the site whose stereo is being interpreted.
\item \textbf{Ring digits at tetrahedral centers.}  The ring digit's lexical
      position contributes to the local ligand order at the chiral atom.
\item \textbf{Implicit hydrogens.}  An implicit hydrogen used as a tetrahedral
      ligand must be represented as an occurrence with a precise placement rule
      in the semantics object.
\item \textbf{Unspecified potential sites.}  Omitting these constraints allows
      accidental stereo to be introduced by slash/backslash tokens used for
      nearby sites.
\item \textbf{RDKit cleanup.}  If RDKit sanitization or cleanup changes atom,
      bond, hydrogen, aromatic, or stereo facts, that must happen before the
      snapshot.  The core enumerator should not run cleanup.
\item \textbf{Symmetry.}  Equal rendered strings may arise from different roots
      or automorphic labeled witnesses.  Use image deduplication or a canonical
      witness rule.
\item \textbf{Contextual atom decoding.}  Nonbracketed atom text may not be a
      unary atom relation.  Model the finite star relation rather than relying
      on a post-hoc parser check.
\item \textbf{Ring endpoint bond conflicts.}  If both endpoints of a ring
      closure emit bond information, compatibility is a joint relation on the
      pair of endpoint slots.
\item \textbf{Policy leakage.}  A presentation convention, such as least-free
      ring labels or canonical annotation, is still a constraint.  It should be
      recorded as policy, not hidden in the renderer.
\end{enumerate}

\section{Testing plan}
\label{s:testing}

Testing should be layered.

\subsection{Core import test}

Assert that core modules do not import RDKit.

\begin{lstlisting}
def test_core_has_no_rdkit_imports() -> None:
    for path in CORE_MODULE_PATHS:
        tree = ast.parse(path.read_text())
        for node in ast.walk(tree):
            if isinstance(node, ast.Import):
                for alias in node.names:
                    assert not alias.name.startswith("rdkit")
            if isinstance(node, ast.ImportFrom):
                assert not (node.module or "").startswith("rdkit")
\end{lstlisting}

\subsection{Fact validation tests}

Hand-build tiny fact objects and test validation failures: duplicate atom ids,
invalid bond endpoints, component partitions that omit atoms, stereo sites with
missing occurrences, and reference orders of the wrong length.

\subsection{Skeleton tests}

Use non-stereo graphs where the expected skeleton count is small: one atom, a
path of length two, a triangle, and two disconnected atoms.  Test that every
skeleton covers all atoms and partitions edges into tree and ring edges.

\subsection{Semantics tests}

Test \texttt{local\_tetra\_order}, \texttt{tetra\_value}, and
\texttt{directional\_value} on hand-built skeletons.  Include cases where a ring
endpoint changes the local tetrahedral order and cases where reversing the
written direction of a slash carrier changes its site-relative sign.

\subsection{Policy tests}

Test support-wise maximal, cardinality-maximal, and canonical annotation on small
artificial carrier-scope examples.  These tests should not need chemistry.

\subsection{RDKit audit tests}

After the core tests, use RDKit as an external audit on a small corpus.  The
audit should parse generated strings and compare a freshly extracted fact object
or independent invariants.  It should never be called by
\texttt{enumerate\_support}.

\section{Minimal viable implementation sequence}
\label{s:mvp}

A realistic build sequence is:
\begin{enumerate}[leftmargin=*]
\item Implement \texttt{MoleculeFacts}, \texttt{GraphIndex}, validation, and
      hand-built molecules.  Do not use RDKit yet.
\item Implement non-stereo traversal skeletons and rendering for a restrictive
      atom/bond text policy.
\item Add bounded ring labels and interval-based label reuse.
\item Add atom, bond, and ring-pair decode relations, still with a restrictive
      policy.
\item Add tetrahedral sites with explicit and implicit ligand occurrences.
\item Add directional sites using the general finite relation interface.
      Implement signed parity as a fast path only after relation tests pass.
\item Add potential unspecified sites and no-accidental-stereo constraints.
\item Add annotation policies and rendered-string deduplication.
\item Add the RDKit adapter and audit tests.
\item Only then add broader atom/bond spelling domains and optional SAT/CP
      backends.
\end{enumerate}

This order prevents RDKit parser behavior from becoming a hidden dependency in
the core design.

\section{Weakest assumptions for no post-hoc parse filter}
\label{s:weakest}

No post-hoc parse filter is required under the following assumptions.
\begin{enumerate}[leftmargin=*]
\item The input is a fixed finite graph with fixed atom, bond, component, and
      stereo facts.
\item The presentation policy has finite syntax domains.
\item The parser semantics for the bounded dialect is explicit as finite
      relations over skeletons and syntax slots.
\item The potential-site universe is complete for that semantics and policy.
\item The renderer is a total function on satisfying assignments and emits only
      grammar-valid syntax.
\item Any parser behavior that would otherwise be procedural---aromaticity,
      implicit hydrogens, valence defaults, stereo cleanup, ring-label reuse,
      or slash/backslash interpretation---is either already fixed in the facts
      or represented inside the finite relations.
\end{enumerate}

If one of these assumptions is dropped, the missing behavior has to reappear as a
post-hoc validator, a writer-specific procedure, or an oracle.  Under these
assumptions, those mechanisms are unnecessary.

\begin{thebibliography}{9}

\bibitem{rdkit-book-stereo}
RDKit documentation, \emph{The RDKit Book}, section ``Stereochemistry''.
The documentation describes the scope of RDKit stereochemistry support, including
tetrahedral atomic stereochemistry and cis/trans stereochemistry at double
bonds.
\url{https://www.rdkit.org/docs/RDKit_Book.html}

\bibitem{rdkit-rdchem}
RDKit documentation, \emph{rdkit.Chem.rdchem module}.
The API documents atom and bond methods, including bond stereochemistry accessors
such as \texttt{Bond.GetStereo()} and \texttt{Bond.GetStereoAtoms()}.
\url{https://www.rdkit.org/docs/source/rdkit.Chem.rdchem.html}

\bibitem{rdkit-findpotential}
RDKit documentation, \emph{rdkit.Chem.rdmolops module}.
The documentation for \texttt{FindPotentialStereo()} says that it finds
potential stereo elements and notes that the function is somewhat experimental
and that the API and results may change.
\url{https://www.rdkit.org/docs/source/rdkit.Chem.rdmolops.html}

\bibitem{daylight-smiles}
Daylight Chemical Information Systems, \emph{SMILES Theory Manual}.
The stereochemistry section discusses local chiral ordering and notes that the
order of a ring-closure bond for chirality is determined by the lexical position
of the ring digit on the chiral atom.
\url{https://www.daylight.com/dayhtml/doc/theory/theory.smiles.html}

\bibitem{opensmiles}
OpenSMILES specification, version 1.0, 15 May 2016.
The specification defines the open SMILES grammar and includes ring-closure and
dot-bond constraints relevant to bounded dialect design.
\url{https://opensmiles.org/opensmiles.html}

\end{thebibliography}

\end{document}
```

[1]: https://www.rdkit.org/docs/RDKit_Book.html?utm_source=chatgpt.com "The RDKit Book — The RDKit 2026.03.2 documentation"

